\chapter{Physics Simulation in Gazebo with ROS 2}
\label{ch:gazebo}

\section{Motivation and Software Stack}

The statistical simulator is fast, seeded and exactly reproducible, but it is
not a physics engine and it does not exercise the questions every deployment
must answer: how the pipeline receives live state from a robot simulator, how
it is packaged so that anyone can reproduce it with one command, and how its
decisions can be observed while a world runs. To answer them the pipeline was
implemented in Gazebo: four differential-drive rovers are simulated by Gazebo
Harmonic, their odometry is bridged into ROS~2 Jazzy, and the full pipeline
runs inside a ROS~2 node against that live stream, publishing its decisions
once per second. The chapter reports an integration, not a new experimental
result.

The integration has four layers (Fig.~\ref{fig:gz:stack}). Gazebo Harmonic
simulates rigid-body physics with the ODE engine, the rover models and their
differential-drive controllers, all described in SDF~1.9. A Gazebo--ROS bridge
forwards the four odometry topics into ROS~2 \emph{Odometry} messages, in the
Gazebo-to-ROS direction only. ROS~2 carries them to the planner node, a thin
adapter around the unchanged pipeline: the pipeline is not rewritten as ROS
packages, because it is middleware-agnostic by design.

\begin{figure}[htbp]
\centering
\begin{tikzpicture}[
  box/.style={draw, rounded corners, align=center, font=\small, minimum height=1.0cm, fill=white},
  arr/.style={-{Stealth[length=2.4mm]}, thick},
  lab/.style={font=\scriptsize, align=center},
]
\node[draw, dashed, rounded corners, minimum width=14.2cm, minimum height=4.6cm] at (6.6,-0.1) {};
\node[font=\small\bfseries, anchor=north west] at (-0.4,2.1) {Docker container (ROS~2 Jazzy base image)};
\node[box, minimum width=3.3cm, fill=orange!8] (gz) at (1.2,0.3) {Gazebo Harmonic\\server (headless)\\world + 4 rovers};
\node[box, minimum width=3.0cm, fill=blue!6] (br) at (6.0,0.3) {Gazebo--ROS\\parameter bridge};
\node[box, minimum width=3.4cm, fill=green!7] (pl) at (11.0,0.3) {planner node\\QUBO + GNN +\\trust model};
\draw[arr] (gz) -- node[above, lab] {odometry\\(Gazebo transport)} (br);
\draw[arr] (br) -- node[above, lab] {Odometry\\(ROS~2, 4 topics)} (pl);
\node[box, minimum width=3.4cm] (st) at (11.0,-1.6) {status and assignment\\topics (JSON)};
\draw[arr] (pl) -- (st);
\end{tikzpicture}
\caption{The Gazebo integration inside one container.}
\label{fig:gz:stack}
\end{figure}

\section{The World and the Rover}

The world is deliberately simple, so that every effect can be attributed
(Table~\ref{tab:gz:world}). Four instances of one rover model stand on a line,
the same line-up as the planned physical testbed. Each rover is a $0.9\times
0.65\times0.25$~m, 3~kg body with two wheels of radius $r_w=0.16$~m on
revolute joints, a track of $L_w=0.76$~m, and a differential-drive plugin that
turns velocity commands into wheel motion and publishes odometry at 20~Hz. The
3D engine of Chapter~\ref{ch:engine} uses the same dimensions, so the two
simulators share one rover.

\begin{table}[htbp]
\centering
\caption{The Gazebo world and the rover model.}
\label{tab:gz:world}
\small
\begin{tabularx}{\linewidth}{@{}l X@{}}
\toprule
Element & Setting \\
\midrule
Physics & ODE, step 1~ms, real-time factor 1.0, update rate 1000~Hz \\
Scene and ground & directional sun with shadows; static $30\times30$~m plane \\
Spawn poses $(x,y,\psi)$ & R0 $(-3.0,0.0,0.00)$, R1 $(-1.0,0.4,0.25)$, R2 $(1.0,0.1,-0.20)$, R3 $(3.0,0.5,0.10)$ \\
Body & box $0.9\times0.65\times0.25$~m, 3.0~kg, inertia $(0.2,0.2,0.3)$~kg\,m$^2$ \\
Wheels and joints & radius 0.16~m, width 0.08~m, 0.25~kg; revolute, unlimited \\
Drive plugin & differential drive, separation 0.76~m, odometry at 20~Hz \\
\bottomrule
\end{tabularx}
\end{table}

A differential-drive rover is commanded by a forward speed $v$ and yaw rate
$\omega$; the plugin converts them into wheel rim speeds and back again for
odometry,
\begin{equation}
v_L=v-\frac{\omega L_w}{2},\quad v_R=v+\frac{\omega L_w}{2};\qquad
v=\frac{v_L+v_R}{2},\quad \omega=\frac{v_R-v_L}{L_w},
\label{eq:gz:diffdrive}
\end{equation}
in Gazebo's convention of counter-clockwise positive yaw with the left wheel at
positive $y$; each wheel turns at $v_{L,R}/r_w$. Chapter~\ref{ch:engine} uses
the same relations in its own mirrored convention and integrates them exactly.

\begin{figure}[htbp]
\centering
\begin{tikzpicture}[scale=0.92]
\draw[gray!25, step=1] (-4.6,-1.6) grid (4.6,2.2);
\draw[-{Stealth}] (-4.6,0) -- (4.8,0) node[right, font=\scriptsize] {$x$ (m)};
\draw[-{Stealth}] (0,-1.6) -- (0,2.4) node[above, font=\scriptsize] {$y$ (m)};
\coordinate (r0) at (-3,0); \coordinate (r1) at (-1,0.4); \coordinate (r2) at (1,0.1); \coordinate (r3) at (3,0.5);
\foreach \a/\b in {r0/r1,r0/r2,r1/r2,r1/r3,r2/r3}{\draw[very thick, blue!60!black] (\a) -- (\b);}
\draw[dashed, red!60!black] (r0) to[bend right=25] node[below, font=\scriptsize] {6.0 m, out of range} (r3);
\foreach \p/\n in {r0/0,r1/1,r2/2,r3/3}{\node[circle, draw, thick, fill=white, inner sep=1.5pt, font=\scriptsize] at (\p) {R\n};}
\end{tikzpicture}
\caption{Top view of the Gazebo world: with a 5~m radio range the four rovers
form five candidate links; the end rovers are 6~m apart and cannot link.}
\label{fig:gz:world}
\end{figure}

\section{The Planner Node}

The node subscribes to the four odometry topics, keeps each rover's latest
planar position, and runs one pipeline cycle per one-second timer tick
(Algorithm~\ref{alg:gz:cycle}). Every cycle it publishes two JSON messages: a
status message (cycle, number of rovers with odometry, candidate and kept
edges, pruner mode, solve time, assignment and every rover's trust) and an
assignment message.

\begin{algorithm}[htbp]
\caption{One cycle of the Gazebo planner node}
\label{alg:gz:cycle}
\begin{algorithmic}[1]
\Require latest odometry positions, trust model, previous kept edges
\For{each pair of rovers within the 5~m radio range}
  \State $P\gets1-d/d_{\max}$; link quality $0.85P+0.15$; cost $0.10+0.50(1-P)$
\EndFor
\State build the QUBO: all four rovers required, root 0, degree cap 3, budget 25 units
\State solve by projected annealing, warm-started from the previous kept edges
\If{no feasible solution} keep every candidate link and log a warning \EndIf
\State score rovers for two tasks (criticality 0.7 and 0.4) on the kept graph
\State multiply scores by $(0.5+0.5\,p_i)$ and assign greedily, one task per rover
\State sample execution and delivery; update trust; age, decay loads, drain batteries
\State publish the status and assignment messages
\end{algorithmic}
\end{algorithm}

Four design choices are worth noting. Link quality is a \emph{deterministic}
function of distance, because Gazebo models geometry rather than radio
propagation; repeated runs are then comparable while Gazebo stays responsible
for where the rovers are. A short annealing schedule (temperatures 0.5 to 0.05,
cooling 0.8, two restarts) solves four rovers in about 13~ms, far inside the
cycle. The node starts from the world's spawn poses, so it decides sensibly
while the bridge is still starting, and its status reports how many rovers have
delivered odometry, so a missing bridge is visible at once. If the physical
graph ever becomes infeasible, the node keeps every link and marks the cycle,
so the demonstration never stops silently.

\section{Packaging and Launch}

Everything runs in one Docker container built on the official ROS~2 Jazzy
image, with the Gazebo simulator, the bridge, the message packages and NumPy
added and the project copied in, so the pipeline is importable as ordinary
Python. A compose file defines one service, with the number of cycles passed as
an environment variable. The launch proceeds as in Fig.~\ref{fig:gz:launch}:
before building, it refuses to continue with less than 8~GB of free disk,
because a truncated build of the large base image fails far from the real
cause; after starting the simulator headless, it waits three seconds and, if
the server has died, prints its log and exits. An optional flag starts the
graphical client when a display is available.

\begin{figure}[htbp]
\centering
\begin{tikzpicture}[
  stage/.style={draw, rounded corners, align=center, font=\small, minimum width=2.45cm, minimum height=1.2cm, fill=blue!4},
  arr/.style={-{Stealth[length=2.3mm]}, thick},
]
\node[stage] (a) at (0,0) {1. check\\$\ge8$~GB disk};
\node[stage] (b) at (2.95,0) {2. source ROS~2,\\model path};
\node[stage] (c) at (5.9,0) {3. start\\bridge};
\node[stage] (d) at (8.85,0) {4. start Gazebo\\headless};
\node[stage] (e) at (11.8,0) {5. health check,\\run planner};
\draw[arr] (a) -- (b); \draw[arr] (b) -- (c); \draw[arr] (c) -- (d); \draw[arr] (d) -- (e);
\end{tikzpicture}
\caption{Launch sequence of the container.}
\label{fig:gz:launch}
\end{figure}

\section{Behaviour and Limitations}

Table~\ref{tab:gz:run} shows the node's decisions on the spawn line-up,
obtained by running its cycle code offline with those poses, which are what the
live node receives while the rovers remain where they spawned.

\begin{table}[htbp]
\centering
\caption[Planner cycles on the spawn line-up (all 5 candidate links kept in every cycle, solve 12 to 16~ms).]{Planner cycles on the spawn line-up (all 5 candidate links kept in
every cycle, solve 12 to 16~ms). Trust is shown after each cycle.}
\label{tab:gz:run}
\small
\begin{tabular}{@{}r c c r r r r@{}}
\toprule
Cycle & Critical task & Routine task & $T_0$ & $T_1$ & $T_2$ & $T_3$ \\
\midrule
0--2 & rover 3 & rover 0 & 0.73$\to$0.92 & 0.50 & 0.50 & 0.73$\to$0.92 \\
3 & rover 3 (fails) & rover 0 & 0.95 & 0.50 & 0.50 & 0.50 \\
4 & rover 1 & rover 0 & 0.97 & 0.73 & 0.50 & 0.50 \\
5--11 & rover 1 & rover 0 & 0.99$\to$1.00 & 0.85$\to$1.00 & 0.50 & 0.50 \\
\bottomrule
\end{tabular}
\end{table}

Two observations are instructive. The pruner keeps all five links every cycle:
on a short, healthy line-up every link's utility is far above the sparsity
penalty and neither the caps nor the budget bind, so keeping everything is
optimal; pruning starts to matter when weak or costly links appear, as in the
larger fleets of Chapter~\ref{ch:results}. And trust visibly steers
assignment: rover~3 wins the critical task while trusted, fails in cycle~3,
drops from 0.92 to 0.50, and the task moves to rover~1 from the next cycle.

The integration is honest about its limits. The node consumes odometry but
publishes no velocity commands, so the rovers currently hold their spawn
poses; adding a waypoint follower that uses the controller of
Chapter~\ref{ch:engine} is the next step and would make the graph change from
cycle to cycle. Link quality has no radio model, outcomes are sampled from
assumed reliabilities, and the world is written for four rovers. Even so, the
unchanged pipeline runs against live simulator state inside standard robot
middleware, is reproducible with one command, and exposes every decision on an
inspectable stream.
